% !TeX root = ../papers/nc-p3.tex
\section{曲線グラフの厳密半安定窓と二直線への特殊化}
\label{sec:curve-graded-limits}
\begin{theorem}\label{thm:curve-strictly-semistable}
左右点ファイバーで一次元のグラフを持つ標準有限Frobenius窓とする。
純関係がなく、四テンソルが厳密半安定なら、全次数で標準Koszul左右AS正則である。
従ってこのグラフにおける非正則窓は、純関係を持つ直線配置か、純関係のない安定軌道に限られる。
\end{theorem}
\begin{proof}
純関係がないので前定理の重み計算から半安定性が成立する。
直線を含む場合は定理\ref{thm:curve-graph-pure-exclusion}、
純な平面三次部分曲線を含む場合は定理\ref{thm:plane-cubic-mark-normality}で正則である。
残る窓の任意の $a\times b$ 部分への関係の制限は階数 $a+b-2$ 以上である。
相異なる寸法は定理\ref{thm:curve-graph-pure-frames}による。等しい場合は
$1,4$ は直接、$2,3$ は同定理の完全交叉の議論を使う。
$2\times2$ の階数1は二直線の同型のグラフを与える。
$3\times3$ の階数3は両次数3の純な局所完全交叉を与え、
一方の平面への閉埋込みでその像は平面三次因子になる。どちらも除外した。

次にある点対の接次元が3なら、定理\ref{thm:maximal-tangent-ore}の最初のJacobi計算により
左右の三次元評価核が一致する。前点での三次縮約の末尾の縮約像がこの核なので、
中支持は高々3である。2以下なら両端支持が3なので、中央で得る三行三列の特異な行列の核で
端を縮約すると非零純関係が生じる（中央支持1のときは端を直接縮約する）。
従って三支持が3である。定理\ref{thm:ternary-point-subalgebras}の巡回移送により
この三テンソルの三方向の二文字縮約空間はいずれも基底変換で窓の関係部分空間になるため、
どれにも純関係はない。同定理で全次数の正則性を得る。

これで全点で接次元が2以下、辺のJacobi階数は4以上である。
定理\ref{thm:all-finite-windows-semistable}の等号方向の証明はこの階数、
各 $a\times b$ での下界 $a+b-2$、純関係の不在、左右点ファイバーだけを使う。
実際、最低重みの重複の排除では低い部分の長方形下界と最高部分のSegre回避を使い、
次の四つの関係重みの評価では $(2,2),(3,1),(4,1),(3,2),(3,3)$ とJacobi階数を使う。
これで相補六重みが確定し、同定理の中央差の一致と下降対の証明をそのまま適用できる。
非自明な等号方向は厳密半安定性から存在し、その極限はPBW正則である。
良い特殊化で元の窓も正則になる。
\end{proof}

\begin{theorem}\label{thm:line-polystable-regular-limit}
定理\ref{thm:four-shift-line-pbw}の四移動不変直線とその上の評価階数3を仮定し、
四本の法線関係の階数が2とする。四テンソルは厳密半安定で、
その唯一の多重安定極限は二組の相補的な直線グラフを持つ標準Koszul左右AS正則PBW窓である。
直線外の点ファイバーは任意でよい。各位置の二文字ずつの直和を選ぶと、
まず法線文字がちょうど二つの六種類のブロックだけに特殊化できる。
隣接二位置が法線文字の場合の四ブロックは非退化二次形式二本の積で、
法線文字が交互の二ブロックは各隣接二対への切断でいずれも可逆な四行四列行列となる。
多重安定極限でもこれらの性質を保つ。
\end{theorem}
\begin{proof}
各文字空間を $N_i\oplus T_i$ とする。隣接二接線文字への制限は
反対二法線文字の非退化関係と接線行列式の積なので、
法線文字が二個未満の項はない。法線文字を重み $1$、接線文字を $-1$ にする
四位置の一径数群での極限 $w_0$ は法線二個の部分である。
法線二個が隣接する各ブロックは上述の積そのものである。

各辺の切断の階数は極限で6以下である。両組の同種文字の間の二ブロックは各階数1で、
混合対どうしの $8\times8$ ブロックは階数4以下になる。その行列を
$(NT,TN)$ の順に四分割すると、非対角の両ブロックは上述の非退化な積を
別の二対へ切断したものなので各階数4である。従って混合ブロック全体が階数4で、
その行・列空間は各混合四空間へ同型に射影する。対角の二ブロックも可逆である。
これで両隣接切断の階数6と交互型の可逆性を得た。

従って $w_0$ は定理\ref{thm:split-inverse-regularity}の可逆二平面分解そのものである。
混合四関係の両射影は同型、二つの同種関係は非退化なので、両側の座標直線はいずれも
左右評価階数3のグラフを与える。定理\ref{thm:split-polystable-good-limits}により
その唯一の閉軌道も同じ可逆二平面分解の正則PBW窓である。
同定理の中心化群の証明で使う不変量は、四つの隣接型ブロックを各々二コピー取り、
各二文字位置で交代形式を縮約した四つの値の積である。
各値は二行列式の積の非零定数倍であり、閉軌道でも四ブロックの非退化性を保つ。
この閉軌道は元の四テンソルの閉包に入り、非零なので元の窓も半安定である。
最初の非自明な一径数群が重み零の極限を持つので安定でない。
極限の交代四テンソルの許容スペクトルとJordan分割は
定理\ref{thm:split-bell-locus}、\ref{thm:split-jordan-transport}に従い、
分解付き同型は定理\ref{thm:split-inverse-transporter}の移送集合の共役類で決まる。
\end{proof}

\begin{corollary}\label{cor:curve-two-by-two-equality}
純関係のない一次元点グラフで、四移動不変の直線成分を持つことと、
全位置で最低・最高各二文字の重みを持つ一径数群によるHilbert--Mumford等号があることは同値である。
このような等号方向の正の重み差は全位置で等しい。
\end{corollary}
\begin{proof}
直線からの方向は前定理にある。逆に最低二空間の双対の直線どうしへの制限は各辺で非零であり、
最高二空間の積だけに支持される関係空間は、Segreを避けるので高々一次元である。
定理\ref{thm:all-finite-windows-semistable}の最初の等号計算から重み差は位置ごとに
$s,t,s,t$ と交互になる。最低関係は反対の最高関係としか対合できない。
従って直線の積への制限は一次元である。その双一次式は点ファイバーから非退化で、
四位置の旗が四移動不変の直線グラフを指定する。各法線関係の階数は純関係の不在で2である。
最低を零とした重みの平均の和は $s+t$ である。
前定理の交互二ブロックはどちらも非零で重みは $2s,2t$ であるから、
等号には $2s,2t\ge s+t$ が必要であり $s=t$ となる。
\end{proof}
